\section{Message Route}
\label{sec:cses-message-route}

\problemstatement{Given an undirected graph of $n$ computers and $m$
connections, find a route with the fewest edges from computer $1$ to
computer $n$. Output the number of computers on the route and the route
itself, or \texttt{IMPOSSIBLE}.}

\begin{patternbox}
Fewest edges in an unweighted graph is \textbf{BFS}. Recording each node's
predecessor lets the shortest route be reconstructed by walking parents
back from $n$ to $1$.
\end{patternbox}

\subsection*{BFS with parent pointers}

BFS from node $1$. The first time a node is dequeued-into, it is via a
shortest route, so store the node it came from. When $n$ is reached, follow
$\textit{parent}$ from $n$ back to $1$, then reverse to get the route in
forward order. If $n$ is never visited, the answer is \texttt{IMPOSSIBLE}.

\begin{ideabox}[Parents Turn Distances into Routes]
This is Labyrinth on a general graph rather than a grid: the same BFS,
the same parent-backtrack. On any unweighted graph, ``fewest edges''
plus ``print the path'' is always BFS plus predecessors.
\end{ideabox}

\subsection*{Algorithm}

\begin{enumerate}
  \item Build adjacency lists. BFS from $1$, storing $\textit{parent}[v]$
        on first visit.
  \item If $n$ unvisited, print \texttt{IMPOSSIBLE}.
  \item Else backtrack $n\to\ldots\to1$, reverse, and print the count and
        route.
\end{enumerate}

\subsection*{Dry run}

\dryrun{Edges $1\!-\!2$, $2\!-\!3$, $1\!-\!3$; $n=3$.}

\begin{itemize}
  \item BFS from $1$ reaches $3$ directly (edge $1\!-\!3$), parent
        $[3]=1$.
  \item Route $1\to3$: two computers. (The two-edge path $1\!-\!2\!-\!3$
        is longer and not chosen.)
\end{itemize}

\subsection*{C++ implementation}

\begin{lstlisting}
#include <bits/stdc++.h>
using namespace std;

int main() {
    int n, m;
    cin >> n >> m;
    vector<vector<int>> adj(n + 1);
    for (int i = 0; i < m; i++) {
        int a, b; cin >> a >> b;
        adj[a].push_back(b);
        adj[b].push_back(a);
    }

    vector<int> parent(n + 1, 0);
    vector<bool> vis(n + 1, false);
    queue<int> q;
    q.push(1); vis[1] = true;
    while (!q.empty()) {
        int u = q.front(); q.pop();
        for (int v : adj[u])
            if (!vis[v]) { vis[v] = true; parent[v] = u; q.push(v); }
    }

    if (!vis[n]) { cout << "IMPOSSIBLE\n"; return 0; }
    vector<int> route;
    for (int v = n; v != 0; v = parent[v]) route.push_back(v);
    reverse(route.begin(), route.end());

    cout << route.size() << "\n";
    for (int v : route) cout << v << " ";
    cout << "\n";
    return 0;
}
\end{lstlisting}

\begin{complexitybox}
\textbf{Time Complexity.} $O(n+m)$ -- one BFS. \textbf{Space Complexity.}
$O(n+m)$ for adjacency lists, parents, and the queue.
\end{complexitybox}

\begin{takeawaybox}
Message Route is the graph (not grid) version of shortest-path-with-route:
BFS for fewest edges, parent pointers for the path. Master this pair --
BFS distance and parent reconstruction -- as the reflex for all unweighted
shortest-path queries.
\end{takeawaybox}
